\begin{description}
\item[(a)] 14--15 days old. \hfill\mbox{[5 points]}
\item[(b)] \(23^\circ\)--\(27^\circ\). The Moon moves about \(1^\circ\) over the
two to three hours of the examination, which the range covers. \hfill\mbox{[5 points]}
\item[(c)] 25--26 days. The Moon is about \(25^\circ\) ahead of Saturn, so it
must travel \(360^\circ-25^\circ=335^\circ\). Using a lunar orbital period of
27.3 days gives \((335/360)\times27.3\simeq25.4\) days. The key treats Saturn's
motion as less than \(1^\circ\) over this interval and neglects it.
\hfill\mbox{[10 points]}
\item[(d)] Measure the Moon's altitude. The Moon is approximately full, so the
Sun is approximately opposite it and its altitude is roughly the negative
of the measured value. The key gives these approximate lunar altitudes for
28 September and accepts an error of \(5^\circ\) either side.
\hfill\mbox{[10 points]}

\begin{center}
\begin{tabular}{@{}lrrrr@{}}
\toprule
Time & 21:00 & 22:00 & 23:00 & 24:00\\
Lunar altitude & \(20^\circ\) & \(30^\circ\) & \(40^\circ\) & \(45^\circ\)\\
\bottomrule
\end{tabular}
\end{center}

\item[(e)] The examiner confirms successful targeting. \hfill\mbox{[5 points]}
\item[(f)] Time how long the Moon and the mare each take to cross the field
of view by their diameter. The ratio of those times equals the ratio of the
angular diameters. The Moon's diameter averages \(30'\) and was to be about
\(33'\) that day. Allowing for the mare's irregular shape, accept \(4'\)--\(6'\).
\hfill\mbox{[10 points]}
\end{description}
